\documentclass[11pt,a4paper]{article}
\usepackage[T1]{fontenc}
\usepackage[utf8]{inputenc}
\usepackage{lmodern,amsmath,amssymb,amsthm,mathtools}
\usepackage[margin=25mm]{geometry}
\usepackage{booktabs,longtable,array,enumitem}
\usepackage[hidelinks]{hyperref}
\usepackage{microtype}
\setlength{\parindent}{0pt}
\setlength{\parskip}{5pt plus 1pt}
\setlist{leftmargin=1.6em,itemsep=3pt,topsep=4pt}
\newtheorem{theorem}{Theorem}[section]
\newtheorem{proposition}[theorem]{Proposition}
\newtheorem{lemma}[theorem]{Lemma}
\newtheorem{corollary}[theorem]{Corollary}
\theoremstyle{definition}
\newtheorem{definition}[theorem]{Definition}
\newtheorem{example}[theorem]{Example}
\newtheorem{remark}[theorem]{Remark}
\newcommand{\Set}{\mathbf{Set}}
\newcommand{\Cat}{\mathbf{Cat}}
\newcommand{\Biprof}{\mathbf{Biprof}}
\newcommand{\W}{\mathcal W_R}
\newcommand{\id}{\mathrm{id}}
\newcommand{\one}{\mathbf 1}
\newcommand{\List}{\mathrm{List}}
\newcommand{\inl}{\mathrm{inl}}
\newcommand{\inr}{\mathrm{inr}}
\newcommand{\psint}{\int_{\mathrm{ps}}}
\newcommand{\supp}{\operatorname{supp}}
\newcommand{\status}[1]{\par\textbf{#1.}\ }
\title{Indexed Receipt Interaction Theorem\\[5pt]
\large Independent weighted contexts and polynomial receipt content}
\author{}
\date{1 November 2026}
\begin{document}
\maketitle
\vspace{-1em}

\status{Conclusion}
TARGET B admits the required complete retained-category interaction. Independent
commutative-semiring and finitary polynomial cartesian-monad inputs derive the
crossing, all native cells and endpoint actions, pure balance, and the full
pseudolaw coherence. The construction has a sharp bounded recognition theorem,
strong coefficient and fixed-$T$ comparisons, and conservative one-sided
controls. Ordinary output-level laws alone do not guarantee this retained
pseudolaw. The theorem package is proved below and independently reviewed;
its scope is not the entire expanded research contract.

\section{Purpose and scope}
\label{sec:scope}

This dossier addresses the authorized full proof attempt for TARGET B of the
O1 feasibility decision~\cite{Feasibility}: independently specified normalized
finite weighted contexts and genuinely indexed receipt content over a finitary
polynomial cartesian monad on $\Set$. The required conclusion is a complete
category-valued pseudodistributive law, not only a function on output
distributions or a multiplication on isomorphism classes. Its structural cells
must be invertible where the definition of a pseudolaw requires invertibility;
native grade arrows need not be invertible.

The proof includes native cells, arbitrary pure endpoint maps,
pseudobalancing, all law-coherence equations, bounded recognition of the
specified recipe, coefficient changes, the fixed-$T$ semantics comparison,
and both one-sided controls. A failure of one of those commitments is not
silently repaired by weakening the theorem. The work does not authorize the
remaining research contract, applications, external submission, or publication.
There is no claim of historical novelty.

\subsection{Inherited results and the precise increment}

The inherited manuscript~\cite{Unified} supplies the relative category-valued
ambient, its concrete realization, the all-endpoint-square construction, and
the pseudocoend/comma equivalence. Its mixed-algebra theorem assumes a complete
coherent crossing. The present task is to derive that crossing from independently
checkable inputs and to prove exactly what the retained interface forces.
The existing finite probability/list example is a control, not an additional
instance to count as a new criterion.

The ordinary Set-level affine/commutative-effect mechanism is also an external
baseline, not a novelty claim: see Fujii--Tsai--Montacute--Hasuo,
Theorem~42~\cite{FTMH}. An ordinary law obtained after forgetting witnesses does
not on its own produce a pseudolaw with retained witness categories. The
relationship with that baseline must therefore distinguish output-level
equations, native-cell lifts, and invertibility of structural comparisons.

\subsection{Relation to the earlier O1 obligation}

The theorem concerns a parameterized positive family, not arbitrary
contextads or contentads. Its content fibres $L_B=\Set/(TB)$ vary with the
interface; this is not TARGET A's constant category of affine containers.
TARGET A remains a separate control, and its affine-cell recognition does not
automatically apply to the label-preserving slice cells of TARGET B.
Even $T=\id$ leaves receipt grades $\Set/B$ and is not a trivialization of
content. Opposite-side trivialization means replacing the entire opposite
relative algebra by the identity; inclusion by the unit grade inside a fixed
mixture is a different statement requiring its own full-faithfulness proof.

The original display/codisplay pullback/pushout programme remains an open
comparison and obstruction branch. It is neither proved here nor silently
identified with the polynomial receipt construction. The successful scoped
proof establishes the authorized structural target; it does not by itself
establish the contrasting interpreted O6 application
families, the no-recovery requirement, a language, QBS comparisons, external
monoidal structure, or the full contract's integrated acceptance conditions.

\section{Independent entry data}
\label{sec:inputs}

Fix nested universes large enough to form the categories and relative
pseudocoends below. The pure base is $\Set$ in the smallest chosen universe.
All set maps are pure. Finite witness carriers and finite polynomial position
sets are required, but interface sets, shape sets, and receipt-tag sets may be
arbitrary small sets in that universe.

\begin{definition}[Raw normalized weighted witnesses]
Let $R$ be a commutative semiring with $0\ne1$. An object of $\W$ is a finite
set $X$ with a weight $p:X\to R$ satisfying $\sum_{x\in X}p(x)=1$. A morphism
$a:(X,p)\to(Y,q)$ is a function satisfying
\[
 q(y)=\sum_{x:\,a(x)=y}p(x)\qquad(y\in Y).
\]
Zero-weight points are retained. No quotient by support, output distribution,
or isomorphism class is taken.
\end{definition}

Set
\[
 K_A=\W^A,\qquad e_AP=\coprod_{a\in A}U(P_a),\qquad
 \pi_P:e_AP\longrightarrow A.
\]
For $u:A'\to A$, reindex by $(u^*P)_{a'}=P_{u(a')}$ and use
$e_u(P)(a',x)=(u(a'),x)$. The unit grade $I_A$ is the singleton family, with
discharge $e_AI_A\to A$. For $D\in K_{e_AP}$, accumulation has fibre
\[
 (P\otimes D)_a=\coprod_{x\in U(P_a)}U(D_{a,x}),
 \qquad w_a(x,z)=p_a(x)d_{a,x}(z).
\]
The extension comparison is the canonical reassociation
$e_A(P\otimes D)\cong e_{e_AP}D$.

Independent tensor in $\W$ means finite Cartesian product of the underlying
sets, with product weights. Its empty tensor is the singleton of weight $1$.
This tensor is not asserted to be a categorical product in $\W$. Finite
distributivity proves normalization and the pushforward formula for products
of witness maps, while the corresponding iterated-sum identity proves it for
dependent accumulation. The canonical tuple rebracketings preserve weights
and all coordinate values. These statements use the semiring axioms, not a
mixed law.

\begin{definition}[Polynomial cartesian content input]
Independently of $R$, fix a monad $(T,\eta,\mu)$ on $\Set$ with a chosen
polynomial presentation
\[
 T(B)=\coprod_{s\in S}B^{E_s},\qquad |E_s|<\infty.
\]
Assume that the naturality squares of $\eta$ and $\mu$ are pullbacks.
Polynomial $T$ preserves pullbacks. Thus $T$ is a finitary polynomial
cartesian monad. These assumptions mention neither weighted contexts nor
an interaction law.
\end{definition}

For each $B$, define
\[
 L^T_B=\Set/(TB),\qquad d_B(q:Q\to TB)=B+Q,
 \qquad \nu_Q=[\eta_B,q]:B+Q\longrightarrow TB.
\]
A native content cell $a:q\to q'$ is a map $a:Q\to Q'$ with $q'a=q$;
$d_Ba=\id_B+a$. For $v:B\to B'$ the covariant action is
$v_!q=Tv\,q$ on the same tag set, and $d_v(Q)=v+\id_Q$.
The content unit is the empty grade $U_B=(\varnothing\to TB)$ with
$B\to B+\varnothing$. If $r:Z\to T(B+Q)$, the ordered accumulation is
\[
 Q\diamond Z=
 \bigl(Q+Z\xrightarrow{[q,\,\mu_BT\nu_Q\,r]}TB\bigr).
\]
Its extension comparison $(B+Q)+Z\to B+(Q+Z)$ is coproduct reassociation.
The monad identities and naturality identify the two label maps in every
unit and associativity comparison; those comparisons do not use a context
grade or a weighting.

\subsection{The retained modules and endpoint equations}

Define
\[
 M_R(A,B)=(e_A\downarrow B),\qquad
 N_T(A,B)=(A\downarrow d_B),\qquad
 H_{R,T}(A,B)=(e_A\downarrow d_B).
\]
A mixed object is $(P,Q,f:e_AP\to B+Q)$. Its fixed-endpoint morphisms are
pairs $(\alpha:P\to P',\beta:Q\to Q')$ satisfying
\[
 (\id_B+\beta)f=f'e_A\alpha.
\]
For pure endpoint maps $u:A\to A'$ and $v:B\to B'$, a native square from
$(P,Q,f)$ to $(P',Q',f')$ is precisely
\begin{equation}\label{eq:endpoint-square}
 \alpha:P\longrightarrow u^*P',\qquad
 \beta:v_!Q\longrightarrow Q',\qquad
 (\id_{B'}+\beta)(v+\id_Q)f
 =f'e_u(P')e_A\alpha.
\end{equation}
The context component is weight preserving in every fibre and the content
component satisfies $q'\beta=Tv\,q$. The retained square is not restricted
to invertible $\alpha$ or $\beta$.
If a second square has endpoints $(u',v')$ and components
$(\alpha',\beta')$, their vertical composite has endpoints
$(u'u,v'v)$ and components
\[
 a\longmapsto\alpha'_{u(a)}\alpha_a,\qquad \beta'\beta.
\]
Its label and body equations are the pasted equations of the two squares;
pushforward preservation is closed under composition. Horizontal pasting
is the explicit dependent context map and lifted receipt map proved in
Section~\ref{pc:composition}. Thus neither pasting operation is defined by
passing to isomorphism classes.

The relative tensor uses execution order:
\[
 (X\odot Y)(A,D)=\psint^{B\in\Set}X(A,B)\times Y(B,D).
\]
Its arrows are generated by native product-category arrows and invertible
pure-middle balance arrows, with their naturality, unit, and composition
relations. The inherited theorem gives a module equivalence
$M_R\odot N_T\simeq H_{R,T}$ without inverting native grade maps.
The new crossing must be defined on that generator-and-relation
presentation, not only on a preferred object representative.

\begin{theorem}[Indexed receipt compatibility and bounded recognition]
\label{thm:main}
Fix a commutative semiring $R$ with $0\ne1$ and the raw witness doctrine
above. For recognition, let $T$ be any finitary polynomial monad, without
initially assuming its unit or multiplication cartesian. Within the fixed
occurrence-normalized, raw-witness, $T(E_P)$-label-retaining receipt class
of Definition~\ref{rb:normalform}, a complete pseudodistributive law
\[
 \tau_{R,T}:N_T\odot M_R\Longrightarrow M_R\odot N_T
\]
exists if and only if $\eta$ and $\mu$ are cartesian.
When these conditions hold, the canonical construction retains all native
cells and endpoint squares, admits coherent strong coefficient changes at
fixed $T$, has strong locally fully faithful one-sided inclusions, and has
a strong receipt-forgetting comparison to the fixed-$T$ history algebra.
The latter comparison is generally neither full nor faithful.
This equivalence is not a classification of arbitrary interaction laws,
correlated sampling strategies, support quotients or lax replacements.
\end{theorem}

The complete crossing proof supplies sufficiency, the retained-normal-form
recognition proof supplies necessity, and the preservation section proves
the stated comparison properties. In particular, invertible structural
comparisons are conclusions of the positive construction, not extra input
compatibility assumptions.

\section{The complete indexed crossing law}\label{xc:law}

Write $M$ for the weighted-context module, $N$ for the receipt-content
module, and tensor in execution order.  We construct
$\tau:NM\Rightarrow MN$. In Gambino--Lobbia's Definition~4.1 and Appendix~A~\cite{GambinoLobbia}
put their $S=N,T=M$, their $(s,m)=(j_N,\mu_N)$ and
$(t,n)=(j_M,\mu_M)$. We use the actual Appendix diagrams, not an informal
grouping of the axioms. The proof below checks C1--C10, hence both their
required C1--C8 presentation and the inherited nine-equation presentation.

\subsection{Polynomial leaves and weighted histories}

Fix a polynomial presentation
\[
 TB=\coprod_{s\in S}B^{E_s},\qquad |E_s|<\infty.
\]
Cartesianness of $\eta$ says its selected shape has one position.
Cartesianness of $\mu$ supplies, for every nested shape,
a bijection
\[
 \rho_{s,(s_i)}:E_{m(s,(s_i))}\xrightarrow{\sim}
                 \coprod_{i\in E_s}E_{s_i}.
\]
These are the backwards position maps representing the polynomial
transformations. Equality of polynomial natural transformations is equality
of their shape maps and backwards position maps: evaluate a representable
summand at its set of distinct formal position variables. The monad
associativity and unit equations therefore imply the following fact.

\begin{lemma}[Polynomial leaf coherence]\label{xc:leaves}
For a fixed finite nested $T$-tree, all monad reassociation and unit-removal
routes give the same flattened shape and the same bijection from its
positions to the original addressed leaves.
\end{lemma}

For a weighted family $P$ over $B$, let
$E_P=\coprod_bP_b$ and $\pi_P:E_P\to B$. Define
\[
 \operatorname{Lift}_P(s,(b_i))=\prod_{i\in E_s}P_{b_i}
\]
with product weight. Its underlying set is canonically the fibre of
$T\pi_P:TE_P\to TB$ over $(s,(b_i))$. A native context map acts
coordinatewise. Summing the product weights over a product of preimages
gives the product of the respective pushforward masses. Thus these maps
preserve weights, including zero-weight witnesses. The dependent comparison
\begin{equation}\label{xc:fubini}
 \prod_i\Bigl(\coprod_{x_i\in P_i}V_{i,x_i}\Bigr)
 \cong
 \coprod_{(x_i)\in\prod_iP_i}\prod_iV_{i,x_i}
\end{equation}
has common point weight
$\prod_i p_i(x_i)v_{i,x_i}(y_i)$. Finite semiring distributivity proves
normalization and pushforward preservation; no division or conditioning
operation is involved. These products are sample tensors, not asserted
categorical products in $\W$.

\subsection{Independent content structure and pullback transport}

For $Q=(q:Q\to TB)$ put $D_BQ=B+Q$ and
$\nu_Q=[\eta_B,q]$. A native map $a:Q\to Q'$ satisfies $q'a=q$;
its extension is $\id_B+a$. Pure transport along $u:B\to B'$ labels
the same carrier by $Tu\,q$, with extension $u+\id_Q$. The unit grade
is empty. For $r:Z\to T(B+Q)$, accumulation has carrier $Q+Z$ and label
\[
 [q,\mu_BT(\nu_Q)r].
\]
Its extension comparison is tagged coproduct reassociation. On native
maps accumulation is the corresponding coproduct function. Its label
equation follows from $\nu_{Q'}(\id+a)=\nu_Q$ and naturality of $\mu$.
The monad laws imply the content unit and associativity equations; the
constraints retain the original tagged summands, so their coherence is
literal equality of these reassociation functions. No sampling is used.

For $f:X\to B$, set $Q^f=Q\times_{TB}TX$ and let
$j_f^Q:X+Q^f\to B+Q$ be the evident map. Cartesianness of $\eta$ and
extensivity give
\begin{equation}\label{xc:pb1}
 X+Q^f\cong(B+Q)\times_{TB}TX.
\end{equation}
Applying $T$, pasting the cartesian square of $\mu$, and pulling back along
$r$ yield
\begin{equation}\label{xc:pb2}
 f^*(Q\diamond Z)\cong Q^f\diamond(j_f^Q)^*Z.
\end{equation}
This isomorphism preserves each original $Q$- or $Z$-tag. On $Z$ it is
the unique lift of the fixed nested structure $r(z)$. All lift-composition
and Beck--Chevalley comparisons here are determined by their projections;
their identity, composition and pasting equations follow by pullback
uniqueness. They are derived content-only comparisons, not supplied mixed
law data.

\subsection{The actual pseudobalanced crossing}

Use the inherited normal category $H(A,D)=(e_A\downarrow d_D)$.
On a representative
\[
 [(Q,f:A\to B+Q),(P,g:E_P\to D)]\in(N\odot M)(A,D)
\]
define $V_a=\operatorname{Lift}_P(\nu_Qf(a))$ and
$Q_P=Q\times_{TB}TE_P$, labelled by $Tg\,\operatorname{pr}_2$.
The body $h:E_AV\to D+Q_P$ sends an ordinary branch
$f(a)=\inl b$ with its one sample $x$ to $\inl g(b,x)$, and a receipt
branch $f(a)=\inr q$ with lifted structure $z$ to $\inr(q,z)$.
Equivalently this is \eqref{xc:pb1} at $f=\pi_P$, followed by
$g+\id$. Denote this construction by $\bar\tau$.

On a native content cell $a:Q\to Q'$, its context map is the identity
on the same $T$-label and its receipt map is $(q,z)\mapsto(aq,z)$.
On a native context cell $b:P\to P'$, its context map applies $b$
coordinatewise, and its receipt map is
$(q,z)\mapsto(q,T(eb)z)$. The native body equation gives equality of
the pushed output labels. The body square holds in both tagged cases;
product pushforward proves the context maps preserve weights. Identity,
composition and commutation of these cell actions follow from function
composition and functoriality of $T$. No native arrows are inverted.

For a pure middle map $u:B\to B'$, the balanced representatives use
$u_!Q,P'$ on one side and $Q,u^*P'$ on the other. Their context fibres
are the same addressed products. Writing $E_u=B\times_{B'}E_{P'}$,
pullback preservation by $T$ gives the receipt isomorphism
\begin{equation}\label{xc:balance}
 Q\times_{TB}TE_u\xrightarrow{\sim}Q\times_{TB'}TE_{P'},
 \qquad(q,z)\longmapsto(q,T(e_u)z).
\end{equation}
This preserves the output labels and bodies, including every unused
receipt. It is an actual pullback isomorphism, not an identity on receipt
grades. Naturality in native maps, and the balance identity and composition
relations, follow by checking its two projections. Hence $\bar\tau$
descends through the generators and relations of the pseudocoend.

Source endpoint maps reindex $V$. Target endpoint maps $v:D\to D'$
replace $g$ by $vg$ and the label by $Tv\,Tg$, exactly the content
pushforward. Their compatibility with balance is again equality of
pullback projections. Thus $\bar\tau$ is a pseudonatural module map.

Every one-sided arrow is generated by its generic arrow and a pure map:
$(P,g)$ is $(P,\id)$ followed by $g$, while $(Q,f)$ is $f$ followed by
$(Q,\id)$. Adjacent pure maps are absorbed by the actual context pullbacks
or content pushforwards. The absorption coherence is the balance and
one-sided action coherence just verified. Consequently the generic
components below determine all components by whiskering and transport;
this does not require intervening bodies to have constant arity.

Let $\Phi:M\odot N\simeq H$ be the inherited global comma equivalence
and $\Psi$ its specified quasi-inverse with triangle constraints. Put
$\tau=\Psi\bar\tau$. We check modifications after $\Phi$, which is
fully faithful and reflects their equality and invertibility. Iterated
tensors are evaluated on their common flattened pseudobalanced presentation.

We use the inherited Gray-presentation convention for the tricategory of
category-valued biprofunctors. Choose its coherent equivalences to the
concrete pseudocoend model, and the unit/counit triangle constraints for
$\Phi\dashv\Psi$ as an adjoint equivalence. Transformations and
modifications below are transported with these constraints; no literal
associativity of relative tensors is assumed. On each iterated tensor, the
pseudobalanced universal property determines a component and a modification
from its values on representatives, native arrows and balance generators.
The preceding verification therefore defines the transported cells, not
merely objectwise functions. In the Gray presentation the standard law
axioms apply; returning along the coherent equivalence inserts precisely
the tensor associators, unitors and interchangers checked in the history
argument below. This is the same inherited transport used for the mixed
algebra theorem, rather than an additional strictification hypothesis.

\subsection{Four derived invertible modifications}

\textbf{Content unit.} For $Q=\varnothing$, its lifted receipt set is empty
and $\operatorname{Lift}_P(\eta b)\cong P_b$ by the one-position unit.
This yields $\tau(j_NM)\cong Mj_N$.

\textbf{Context unit.} If each $P_b$ is singleton, all tuple coordinates
are forced, including nullary shapes; $\pi_P$ is an isomorphism and
$Q_P\cong Q$. Deleting the forced coordinates gives
$\tau(Nj_M)\cong j_MN$.

\textbf{Content multiplication.} For nested content $Q,Z$, use
\eqref{xc:pb2} at $\pi_P$. On every receipt retain its original $Q$ or
$Z$ tag; on $Z$, replace a flat lift of $\mu T\nu_Qr(z)$ by its unique
nested lift. Lemma~\ref{xc:leaves} identifies flat and nested sample
positions. The same factors occur in the weight product. This gives
\[
 \tau(\mu_NM)\cong(M\mu_N)(\tau N)(N\tau).
\]

\textbf{Context multiplication.} For $P$ followed by $V$, split each
occurrence sample $(x_i,y_i)$ into its tuple $(x_i)$ and dependent tuple
$(y_i)$. On contexts this is \eqref{xc:fubini}; on receipts it is the
same map in the pullback-composition identification along
$E_{P\cdot V}\cong E_V\to E_P\to B$. It preserves the original tag,
intermediate $TE_P$-label and all values. This gives
\[
 \tau(N\mu_M)\cong(\mu_MN)(M\tau)(\tau M).
\]

These maps have the stated inverses, preserve all raw witnesses and weights,
and satisfy the body equations. Their naturality follows from product
functoriality, pullback projections and dependent pushforward. The formulas
apply to all grade elements, not only those reached by a body.

\subsection{A faithful normal form for the coherence check}

A context history records its input, every ancestor sample, every intermediate
$T$-shape selected by the original bodies, and a sample at each valid
addressed leaf and context stage. Later fibres may depend on earlier values;
these are dependent valid-history sets, not products over a falsely fixed
index set. A receipt history additionally retains the original content-stage
number and original receipt tag, its nested label, and every descendant
sample choice. It ranges over \emph{all} tags in the grade sums, including
unused alternatives.

Each bracketing occurring in C1--C10 expands bijectively to the corresponding
common history set. Dependent-sum reassociation retains ancestor values;
product/sum exchange is \eqref{xc:fubini}; content splitting is the unique
lift in \eqref{xc:pb2} with the leaf bijection of
Lemma~\ref{xc:leaves}; coproduct reassociation retains stage tags; units
remove only forced coordinates or empty summands; pure balance retains the
two projections in \eqref{xc:balance}. These observations give an explicit
recursive inverse to expansion for each bracketing. Intermediate shapes
are reconstructed from the retained original labels and ancestor samples,
so flattening does not lose variable-length or nullary provenance.

Every comparison above commutes with expansion. Equality of functions on
both sorts of histories is a faithful equality test: a comma cell is exactly
its context-family functions and its receipt-set function. This includes
zero-weight points; it is not equality only almost surely or equality of
output distributions. All native maps and endpoint transports act by their
actual sample, tag and label functions, so the same test establishes
naturality.

\subsection{The complete coherence inventory}

For clarity, the exact source correspondence is
\[
\begin{array}{c|cccccccccc}
\text{GL Appendix}&C1&C2&C3&C4&C5&C6&C7&C8&C9&C10\\\hline
\text{Marmolejo coh}&4&2&5&3&1&6&9&7&\text{none}&8.
\end{array}
\]
The following checks the actual Appendix boundaries with $S=N,T=M$.
Each calculation concerns both context and receipt histories.

\begin{description}[style=nextline]
\item[C1 content associativity on $NNNM$.]
Retain the original tag in $Q+Z+W$ and all sample values at its original
three-level $T$-tree leaves. The two routes use
$\mu(\mu T)=\mu(T\mu)$, hence have the same shape and leaf bijection by
Lemma~\ref{xc:leaves}. They retain the same root stage and the same product
weight, so their history functions coincide.

\item[C2 content left unit.]
An empty receipt grade introduces no tag; the unit $\eta$ introduces one
position. Both the multiplication-comparison route and the content-unitor
route remove that same empty summand and one-position node, using the
one-coordinate content-unit comparison. Both act identically on remaining
histories.

\item[C3 content multiplication against the context unit.]
All sample fibres are singleton. Both routes on the $NN$ boundary retain
the same original $Q$ or $Z$ tag and remove the unique forced tuple given
by \eqref{xc:pb2}. Context maps are the unique singleton maps. This checks
the unit of $M$ against multiplication of $N$.

\item[C4 the two units.]
With empty receipt grades and singleton context, each input has one history
and there are no receipts. Both pastings give that history and the same
pure body. The pure-action interchanger has this same action.

\item[C5 the two multiplication comparisons on $NNMM$.]
A record retains the original $Q/Z$ tag, its original nested $T$-structure,
and at every original leaf the dependent pair $(x_i,y_i)$. One route
flattens leaf blocks before splitting the pairs; the other splits inside
blocks before flattening. Both give the same pair at each original address,
retain the same root tag and have weight
$\prod_i p_i(x_i)v_{i,x_i}(y_i)$. This verifies the interaction-of-multiplications
pasting, including the interchange of independently applied component maps;
it is not inferred merely from the four modification boundaries.

\item[C6 content unit against context multiplication.]
The content unit has exactly one occurrence. Both routes send the dependent
pair $(x,y)$ to itself, with no receipt tag and weight $p(x)v_x(y)$.
The one-position Fubini comparison is this same pair identification.

\item[C7 context associativity on $NMMM$.]
At every original leaf retain the dependent triple $(x_i,y_i,z_i)$.
The two routes expand $((x,y),z)$ and $(x,(y,z))$ into the same triple,
retaining the original receipt tag and intermediate shapes. The common
weight is
$\prod_i p_i(x_i)v_{i,x_i}(y_i)w_{i,x_i,y_i}(z_i)$.
Thus both history functions coincide even when later arities depend on
previous samples.

\item[C8 context right unit.]
The inserted singleton stage contributes only forced coordinates. Removing
it before or after crossing retains exactly the same history and receipt
decoration. The removed weight factors are $1$.

\item[C9 content right unit.]
The trailing empty receipt summand and its $\eta$-node are removed by both
routes. The equation $\mu(T\eta)=\id$ identifies their shape and position
bijections. Thus the redundant right-content-unit axiom holds directly.

\item[C10 context left unit.]
An initial singleton stage is inserted. Splitting and deleting it, or
removing it before crossing, leaves the same samples at every leaf and the
same original tags. All extra factors are $1$. Thus this redundant GL axiom,
which is Marmolejo's equation (coh~8), also holds directly.
\end{description}

Ambient tensor associators simply flatten the list of middle representatives
and balance generators; their history action is the displayed reassociation.
Gray interchangers compare application of natural maps in independent slots;
their actions are the equal coordinate/tag functions just used in C5.
The history proof therefore includes the ambient constraints of the stated
full law definition. Faithfulness of the history projection and of $\Phi$
reflects each equality back to the original relative tensor.

\begin{theorem}[Derived retained-witness law]\label{xc:theorem}
The independent commutative-semiring and finitary polynomial cartesian-monad
data determine the crossing $\tau$, the four invertible modifications and
all full pseudodistributive-law coherence equations at the stated scope.
The inherited mixed-algebra theorem consequently supplies a relative
pseudopromonad on the actual comma categories $(e_A\downarrow d_B)$,
retaining native context maps, strict $T$-label-preserving receipt maps and
all changing-endpoint squares.
\end{theorem}

This conclusion neither supplies arbitrary contentads nor permits affine
position-deleting receipt cells. Recognition, coefficient comparisons,
one-sided controls and the fixed-$T$ forgetting comparison require their
separate companion arguments. No support quotient, disintegration or
higher-order semantic closure is used in the proof.

\section{Bounded recognition and a sharp retained-law boundary}\label{rb:section}

This section distinguishes recognition of the canonical content recipe from
recognition of a complete retained interaction in a specified normal-form
class. It does not classify arbitrary mixed laws. The necessity result uses
raw witness sets, including zero-weight points, as stipulated throughout.

\subsection{Polynomial coordinates and cartesianness}\label{rb:coordinates}

Fix a finitary polynomial presentation
\[
 TX=\coprod_{s\in S}X^{E_s},\qquad |E_s|<\infty.
\]
A natural transformation between polynomial functors consists of a shape
map and, at each source shape, a backwards map from target positions to
source positions. This follows by evaluating a representable summand at
its universal element $\id:E_s\to E_s$. Such a transformation is cartesian
if and only if every backwards position map is bijective. Sufficiency is
the unique coordinate-lifting property. For necessity, apply the naturality
square to $X\to\one$: in a fixed source-shape fibre its comparison is the
coordinate map from $X$-labels on source positions to $X$-labels on target
positions. For $X=\{0,1\}$ this is bijective exactly when the backwards map
is bijective. The shape map itself need not be injective.

Write the unit shape as $u$, with its unique position map $E_u\to\one$.
A nested shape consists of $s\in S$ and inner shapes $t_i$ for $i\in E_s$.
Put
\[
 L=\coprod_{i\in E_s}E_{t_i}.
\]
Multiplication chooses an output shape $m=m(s,(t_i))$ and a map
$\kappa:E_m\to L$. Consequently $\eta$ is cartesian precisely when
$|E_u|=1$, and $\mu$ is cartesian precisely when every $\kappa$ is bijective.
A polynomial functor preserves pullbacks; its monad unit and multiplication
need not be cartesian.

\subsection{The content-only recipe characterization}\label{rb:recipe}

For $q:Q\to TB$, let $\nu_Q=[\eta_B,q]:B+Q\to TB$ as above. For
$f:X\to B$, put $Q^f=Q\times_{TB}TX$ and
$j_f:X+Q^f\to B+Q$. The occurrence-display square is
\[
\begin{array}{ccc}
X+Q^f&\xrightarrow{j_f}&B+Q\\
\nu_{Q^f}\downarrow&&\downarrow\nu_Q\\
TX&\xrightarrow{Tf}&TB.
\end{array}
\]
For $r:Z\to T(B+Q)$, the canonical substitution comparison is the
tag-preserving map
\[
 Q^f+Z^{j_f}\longrightarrow(Q+Z)^f.
\]
The accumulated $Z$-label is $\mu_BT(\nu_Q)r$.

\begin{proposition}[Canonical recipe recognition]\label{rb:recipe-iff}
Assume $T$ preserves pullbacks. The occurrence-display square is a pullback
for every $f,Q$, and the canonical tag-preserving substitution comparison
is bijective for every $f,Q,Z$, if and only if $\eta$ and $\mu$ are
cartesian.
\end{proposition}
\begin{proof}
For sufficiency, cartesianness of $\eta$ and disjoint coproducts identify
$X+Q^f$ with $(B+Q)\times_{TB}TX$. Apply $T$, paste the cartesian
naturality square of $\mu$, and pull back along $r$. This gives the
comparison on the $Z$-summands; the $Q$-summands are already $Q^f$.
These maps are determined by pullback projections, independently of weights.

For necessity, taking $Q=\varnothing$ in the display square gives exactly
the naturality square of $\eta$. For multiplication take
\[
 Q=TB,\quad q=\id_{TB},\quad Z=TTB,\quad
 r=T(\inr):TTB\longrightarrow T(B+TB).
\]
Then $Q^f=TX$. Applying pullback-preserving $T$ to the square of coproduct
inclusions identifies $Z^{j_f}$ with $TTX$. On the $Z$-summand the
canonical comparison is
\[
 TTX\longrightarrow TTB\times_{TB}TX,\qquad
 w\longmapsto(TTf(w),\mu_X(w)).
\]
Its bijectivity is cartesianness of $\mu$ at $f$. Restriction to this
summand is justified because the comparison is canonical and tag-preserving;
an unspecified isomorphism of the total coproducts would not suffice.
\end{proof}

\subsection{The retained occurrence normal-form class}\label{rb:class}

Fix a commutative semiring $R$ with $0\ne1$, and the raw finite normalized
sample category $\W$. Native maps are deterministic pushforward-preserving
functions, so invertible native maps are underlying bijections. A native
map of context families over $A$ fixes the original index in $A$.

The following definition states the precise class to which the converse
applies. It fixes the interaction's sampling and observable label structure,
but does not assume that its coherence maps are the canonical ones.

\begin{definition}[Retained occurrence-normalized class]\label{rb:normalform}
A crossing belongs to this class if, after the inherited relative-tensor
to mixed-comma equivalence, it has the following generic normal forms.
\begin{enumerate}
\item For $q:Q\to TB$, a context $P$ over $B$, and
$y\in B+Q$ with $\nu_Q(y)=(s,b:E_s\to B)$, its context fibre is the
independent product $\prod_{i\in E_s}P_{b(i)}$, with its prescribed raw
carrier and product weights.
\item Write $\pi:E_P\to B$ for the total context projection. Output
receipts have the prescribed grade
$\widehat Q=Q\times_{TB}TE_P\to TE_P$, or a specified isomorphic
presentation. On a receipt input, the output's $TE_P$-label is the same
$T$-shape with every occurrence decorated by its chosen sample. Native
output-grade maps preserve that label.
\item The unit and content-multiplication comparisons required by a
pseudodistributive law are invertible retained cells at the fixed module
endpoints. They may be arbitrary such cells and may rename receipts
subject to the native label equations. Their context carriers and body
labels remain those in the preceding two clauses, up to specified
occurrence-preserving isomorphisms.
\end{enumerate}
\end{definition}

The tensor/comma equivalence is fully faithful and reflects invertibility.
Thus an invertible comparison becomes a comma isomorphism, with a native
context bijection in each fixed input fibre and a receipt component
preserving the $TE_P$-label. A different intermediate tensor factorization
does not evade these conditions. This definition excludes differently
normalized or correlated sampling, witness quotients, and lax laws.

\begin{theorem}[Necessary cartesianness for retained invertible coherence]
\label{rb:necessity}
Let $T$ be a finitary polynomial monad, and let $R$ be a commutative
semiring with $0\ne1$. If a crossing in the retained occurrence-normalized
class has the invertible unit and content-multiplication comparison
components required by a full pseudodistributive law, then $\eta$ and
$\mu$ are cartesian.
\end{theorem}
\begin{proof}
For the unit choose $B=\one$ and a two-point raw normalized sample $E$.
Such an object exists over every allowed $R$: assign its points weights
$(1,0)$. If the unit shape has $n$ positions, crossing the content unit
gives context carrier $E^n$, while the other side of the unit comparison
has carrier $E$. An invertible native context component over the unique
input forces $2^n=2$, hence $n=1$. Therefore $\eta$ is cartesian.

For multiplication fix any nested shape $(s,(t_i))$ and its backwards
position map $\kappa:E_m\to L$. Again use $B=\one$ and the two-point
raw context $E$. Let $Q$ have one distinct receipt $q_i$ for each outer
position $i\in E_s$, labelled by the inner shape $t_i$ with all $B$-labels
equal to the unique point. Let $Z=\{z\}$ and set
\[
 r(z)=(s,i\mapsto\inr(q_i))\in T(B+Q).
\]
Use the generic representative with input $A=(B+Q)+Z$ and terminal
context body $\id:E\to E$. Orient the invertible multiplication cell from
the nested-crossing normal form to the flat-crossing normal form, inverting
it if necessary. At the original input $\inr(z)$, its context component is
a bijection
\[
 \alpha:E^L\xrightarrow{\sim}E^{E_m}.
\]
For $w\in E^L$, the nested output's receipt label is
$(m,w\kappa)$, because accumulation applies $\mu$ to the nested decorated
$T$-structure. The flat output's label is $(m,\alpha(w))$. Apply the
receipt label map into $TE$ to the comma-cell body equation. Since native
receipt maps preserve this label, it forces
\[
 \alpha(w)=w\kappa\qquad\text{for every raw tuple }w.
\]
Thus the context component is forced to be the coordinate map, although
it was not assumed canonical. Receipt renaming cannot alter a $TE$-label.
Nor can the context component move to a different original receipt input,
because it is a map of families over the fixed $A$.

If $\kappa$ is not injective, two target coordinates must agree, so this
coordinate map is not surjective on the two-point set. If $\kappa$ is not
surjective, changing an unused source coordinate gives distinct tuples
with the same image, so the map is not injective. Bijectivity of $\alpha$
therefore forces $\kappa$ bijective. Since the nested shape was arbitrary,
$\mu$ is cartesian.
\end{proof}

Nullary and repeated inner shapes are included: distinct receipt tags retain
their positions even when their $T$-labels coincide, and empty products
have their singleton carrier. The proof uses all raw tuples, including
zero-weight ones. If those points are removed or identified, the uniform
statement for every $R$ does not follow. A full-support two-point probability
object gives the same proof without zero-weight witnesses. This is an
\emph{invertible raw-cell} obstruction, unlike the affine naturality
criterion for noninvertible weight-preserving maps, which requires a
nonzero off-diagonal weight.

\begin{corollary}[Exact bounded criterion]\label{rb:iff}
Fix the receipt recipe, raw normalized $R$-witnesses, and the retained
occurrence-normalized class of Definition~\ref{rb:normalform}. For a
finitary polynomial monad $T$, a complete retained pseudodistributive law
of the specified form exists if and only if $T$ has cartesian unit and
multiplication.
\end{corollary}
\begin{proof}
The complete crossing construction of Section~\ref{xc:law} proves
sufficiency from those independently specified hypotheses.
Theorem~\ref{rb:necessity} proves necessity. Polynomiality supplies
pullback preservation. No assertion is made about interactions outside
the specified normal-form class.
\end{proof}

\subsection{An ordinary law without the retained pseudolaw}\label{rb:prefix}

Let $T_{\le2}X$ be the nonempty lists of length at most two:
\[
 T_{\le2}X=X+X^2.
\]
Write $\rho$ for keeping the first two entries of a nonempty list. The unit
forms a singleton; multiplication concatenates a nonempty list of one or
two nonempty short lists and applies $\rho$. Naturality is immediate.
The unit identities hold because each list already has length at most two.
Both associativity composites return the first two leaves of the same
nonempty three-level word. Truncating an inner block cannot remove either
of those first two leaves, and discarding outer blocks after the second
is safe because each block is nonempty. This proves the monad laws.
Equivalently, its associative operation is concatenation followed by
$\rho$, satisfying $xyz=xy$, and its free algebra on $X$ is $X+X^2$.

The functor is finitary polynomial and its unit is cartesian. Its
multiplication is not cartesian: at two binary inner blocks the backwards
position map includes the first two positions into the four leaves.
For example, under $f:\{0,1\}\to\one$, the distinct nested lists
\[
 [[0,0],[0,0]],\qquad [[0,0],[1,1]]
\]
have the same truncated flattening $[0,0]$ and the same $TTf$-image
$[[*,*],[*,*]]$.

Let $D$ be the ordinary finitely supported probability monad, or normalized
finite distributions over a commutative semiring. Define
$\lambda:T_{\le2}D\Rightarrow DT_{\le2}$ by independently sampling the
one or two positions and retaining their list shape. This is a complete
ordinary distributive law, as all four equations show:
\begin{enumerate}
\item The $T_{\le2}$-unit equation says singleton traversal equals
$D$ applied to the singleton map.
\item The $D$-unit equation says sampling Dirac inputs produces the
Dirac distribution at the original short list.
\item The $D$-multiplication equation is finite distributivity: products
of finite mixtures equal the corresponding mixture of products. In
coordinates both sides sum the same products of mixture and sample
weights.
\item The $T_{\le2}$-multiplication equation follows by marginalization.
Flattening and truncating retains the first two occurrence factors.
Sampling all nested occurrences and then truncating sums each unused
normalized factor to $1$, giving exactly the distribution obtained by
sampling only the retained factors. Nonempty blocks ensure both paths
retain the same two occurrences.
\end{enumerate}
This ordinary mechanism is standard affine/commutative-effect mathematics;
it is not claimed as a new distributive law.

In the retained receipt recipe, take $B=\one$, two inner receipt tags each
labelled $[*,*]$, and one outer tag labelled by that pair of inner receipts.
For a fair-bit context, the nested multiplication route has four independent
samples and sixteen raw points; the flat route has two samples and four
raw points. An invertible native context component cannot exist. The
canonical projection to the first two samples is weight-preserving but
noninvertible. A corresponding receipt map may forget the dropped
witnesses, so a local marginalizing comparison exists; this statement is
not a proof of a complete lax-law repair.

Thus an ordinary output-distribution law can survive when the full
witness-retaining pseudolaw fails. This is outside the positive criterion
precisely through noncartesian multiplication, not a counterexample to it.

\begin{remark}[Reader caution]\label{rb:reader}
The reader polynomial monad $X^{\{0,1\}}$ also gives sixteen versus four
points at multiplication and four versus two at the unit. However,
independent traversal does not define a complete ordinary probability law
for that monad: its unit duplicates a distribution, giving independent
draws on one side and the diagonal correlated distribution on the other.
The reader example must not be used as the ordinary-positive,
retained-negative contrast; $T_{\le2}$ supplies that contrast.
\end{remark}

\subsection{Scope of the recognition result}\label{rb:scope}

Cartesianness of the supplied polynomial monad and the semiring witness
doctrine are independent input conditions. Theorem~\ref{rb:necessity}
connects these conditions to the full interaction's retained interface,
while Proposition~\ref{rb:recipe-iff} independently recognizes the
canonical content-only reassignment recipe. Neither statement is a
classification of arbitrary contextads, arbitrary indexed contentads,
support quotients, correlated alternatives, or lax laws. The original
general display/codisplay problem is outside this theorem. Historical
novelty is not implied by the proof or by the more elaborate ambient.


\section{Preservation and conservative controls}\label{pc:section}

This section uses the canonical mixed composition and addressed-history
associator constructed above. It proves the parameter comparisons and the
one-sided and fixed-$T$ controls separately from the existence of the mixed
law. Throughout, $T$ is fixed. Coefficient change and receipt forgetting
have different information-loss properties: the former retains all native
cell functions, while the latter can identify distinct receipt cells.

\subsection{An explicit composite and its native squares}\label{pc:composition}

Write a mixed arrow $F:A\to B$ as
\[
 F=(P,q:Q\to TB,f:E_P\to B+Q),\qquad
 \nu_Q=[\eta_B,q]:B+Q\to TB.
\]
A square to $F'=(P',q':Q'\to TB',f')$ over
$u:A\to A'$ and $v:B\to B'$ consists of
\[
 \alpha_a:P_a\to P'_{u(a)},\qquad
 \beta:Q\to Q',\qquad q'\beta=Tv\,q,\qquad
 (v+\beta)f=f'\bar\alpha,
\]
where $\bar\alpha(a,x)=(u(a),\alpha_a(x))$ and every $\alpha_a$
preserves the specified weighted pushforward. Distinct functions $\beta$
remain distinct cells even when the body never uses the affected receipt.

For $G=(R,s:S\to TC,g:E_R\to C+S):B\to C$, put
\[
 \pi=\pi_R:E_R\to B,\qquad Q_R=Q\times_{TB}T(E_R),\qquad
 \Gamma_g=\mu_C T(\nu_S)Tg:T(E_R)\to TC.
\]
For $\nu_Q f(a,x)=(\sigma,(b_i)_{i\in E_\sigma})$, the composite
$F;G$ has context and receipt grade
\begin{align*}
 H_a&=\coprod_{x\in P_a}\prod_{i\in E_\sigma}R_{b_i},\\
 U&=S+Q_R,\qquad
 \ell_U|_S=s,\quad \ell_U(q_0,t)=\Gamma_g(t).
\end{align*}
The weight at $(x,(r_i))$ is $p_a(x)\prod_i r_{b_i}(r_i)$.
If $f(a,x)=\inl b$, the cartesian unit identifies its single position
with a sample $r\in R_b$, and the composite body is $g(b,r)$, included
from $C+S$ into $C+(S+Q_R)$. If $f(a,x)=\inr q_0$, the sample tuple
determines $t\in T(E_R)$ over $q(q_0)$ and the body returns the receipt
$(q_0,t)$. Thus, in both branches,
\begin{equation}\label{pc:label-composite}
 \nu_U h(a,x,(r_i))=\Gamma_g(t),
\end{equation}
where the ordinary branch has $t=\eta_{E_R}(b,r)$. The ordinary case
uses naturality and the unit law of $\eta$; the receipt case is the
definition of $\ell_U$. Notice the order $S+Q_R$: later receipts precede
lifted earlier receipts. All receipts are included, not only those reached
by the body.

Let $(\alpha,\beta):F\Rightarrow F'$ have endpoints $(u,v)$, and
$(\gamma,\delta):G\Rightarrow G'$ have endpoints $(v,w)$. Set
$\bar\gamma(b,r)=(v(b),\gamma_b(r))$. The map of lifted receipts is
\begin{equation}\label{pc:lifted-cell}
 \beta_R(q_0,t)=(\beta(q_0),T\bar\gamma(t)).
\end{equation}
It lands in the pullback because
$T\pi'\,T\bar\gamma=Tv\,T\pi$ and $q'\beta=Tv\,q$.
The composite receipt cell is $\delta+\beta_R$. Its label equation follows
from
\begin{equation}\label{pc:label-naturality}
 \Gamma_{g'}T\bar\gamma=Tw\,\Gamma_g,
\end{equation}
using $g'\bar\gamma=(w+\delta)g$,
$\nu_{S'}(w+\delta)=Tw\,\nu_S$, and naturality of $\mu$.

The composite context cell sends
\[
 (x,(r_i))\longmapsto(\alpha(x),(\gamma_{b_i}(r_i))).
\]
The positions agree canonically because
$\nu_{Q'}f'\bar\alpha=Tv\,\nu_Qf$. For a fixed target $x'$ and
tuple $(r'_i)$, summing over the preimages of each $\gamma_{b_i}$ gives
the product of the target coordinate weights. That product is independent
of the remaining choice of $x$ over $x'$, since the target labels are fixed
by $f'(u(a),x')$. Summing over $\alpha(x)=x'$ then gives the target weight
of $x'$. Finite distributivity proves the composite pushforward equation.
Equations~\eqref{pc:label-composite}--\eqref{pc:label-naturality}, separately
on ordinary and receipt outputs, prove the body square. This calculation
includes arbitrary noninvertible cells and changing endpoints.

\subsection{Change of coefficient semiring}\label{pc:coefficients}

\begin{proposition}\label{pc:coefficient-comparison}
A unital semiring homomorphism $\varphi:R\to R'$ induces a pure-compatible
strong comparison
\[
 C_\varphi:H_{R,T}\longrightarrow H_{R',T}.
\]
It is locally faithful. It is locally full if $\varphi$ is injective, but
need not be locally full otherwise. Identity and composite coefficient
changes are coherent.
\end{proposition}
\begin{proof}
On weighted sets put
$C_\varphi(X,p)=(X,\varphi p)$ and retain every function. Preservation of
finite sums and $1$ proves normalization and preservation of pushforward
equations. Products and dependent accumulation have the same carriers,
with equal weights on the two comparison routes:
\[
 \varphi\Bigl(\prod_i p_i\Bigr)=\prod_i\varphi(p_i),\qquad
 \varphi(p(x)r_x(y))=\varphi(p(x))\varphi(r_x(y)).
\]
Define
\[
 C_\varphi(P,Q,f)=(C_\varphi P,Q,f),\qquad
 C_\varphi(\alpha,\beta)=(\alpha,\beta).
\]
The receipt and body equations, including arbitrary pure endpoints, are
unchanged. In the composite of Section~\ref{pc:composition}, $Q_R$ depends
on the witness carrier, not its weights, so it is unchanged. The two context
carriers are identical and their weights agree by the displayed equations.
The compositor is the identity on the addressed context and receipt
carriers, or the canonical identity-carrier comparison under weak choices.
The unit comparison fixes the singleton of weight $1$.

Naturality follows from the explicit context cell and
\eqref{pc:lifted-cell}. Every route in the strong comparison pentagon and
triangle has the same underlying context and receipt functions: coefficient
change affects neither rebracketing nor the pullback comparisons. Equality
of their weights was already checked. Thus these equations hold as native
cells, relative to the canonical addressed-history associator, not merely
after forgetting to output laws. The pure embedding is fixed.

For $\psi:R'\to R''$, the equalities
$C_\psi C_\varphi=C_{\psi\varphi}$ and $C_{\id}=\id$ hold on weights,
carriers, bodies, cell functions, and structural comparisons. This gives
the asserted coherent parameter dependence.

Local faithfulness holds because neither $\alpha$ nor $\beta$ is identified.
If $\varphi$ is injective, an $R'$-pushforward equation between images
reflects to the corresponding equation in $R$, proving local fullness.
For a counterexample without injectivity, use the homomorphism
$\mathbb Q_{\geq0}\to\mathbb B$ sending positive weights to true. A
surjection from a uniform three-point carrier to a uniform two-point carrier
preserves the resulting Boolean full-support weights. No such function
pushes rational weights $1/3$ to weights $1/2$. Constant ordinary bodies
and empty receipt grades turn it into a target native cell with no source
preimage. No equivalence is asserted.
\end{proof}

The positive-to-Boolean comparison is substantive: it forgets numerical
probability while retaining the raw possibility histories, including any
zero-weight points. It is not a representation theorem identifying the
two semantics.

\subsection{The two one-sided embeddings}\label{pc:one-sided}

\begin{proposition}\label{pc:one-sided-prop}
The context-only and receipt-content-only structures admit pure-compatible
strong embeddings into $H_{R,T}$ that are fully faithful on every hom
category and retain exactly their changing-endpoint squares.
\end{proposition}
\begin{proof}
For context-only $(P,k:E_P\to B)$ set
\[
 i_M(P,k)=(P,\varnothing,\inl k).
\]
The receipt component of a cell between images is the unique map
$\varnothing\to\varnothing$. Its remaining equation is exactly
$k=k'E_\alpha$ at fixed endpoints, or $vk=k'\bar\alpha$ at endpoints
$(u,v)$. This proves the hom and square claims.

For composition, $\eta_B(b)$ has exactly one position. Hence the mixed
context is canonically $\coprod_x R_{k(x)}$ with weight
$p(x)r_{k(x)}(y)$, and the receipt grade is empty. The body is the native
context composite. The resulting compositor and unit comparison are
invertible. Their coherence removes the unique unit-occurrence address
at each stage and identifies exactly the same nested sample history;
weights agree by product regrouping. Thus $i_M$ is strong.

For content-only $(Q,f:A\to B+Q)$, take singleton weight-$1$ contexts
$I_A(a)$ and set
\[
 i_N(Q,f)=(I_A,Q,f\pi_{I_A}).
\]
The context map $I_A\to u^*I_{A'}$ is unique, so mixed cells and endpoint
squares between these images are precisely the original content cells
and equations. This proves local full faithfulness and the square claim.

Every product of singleton samples is singleton, including the empty
product. The composite context is therefore canonically $I_A$.
For the later context $I_B$, $\pi:E_{I_B}\to B$ is a bijection, and
\[
 Q\times_{TB}T(E_{I_B})\cong Q
\]
canonically via $T\pi$. The composite receipt grade becomes $S+Q$,
with labels $s$ on $S$ and $\mu_C T\nu_S Tg\,q$ on $Q$.
This is exactly the native content composition: push the earlier $Q$
along $g:B\to C+S$, then accumulate it after $S$. The bodies agree in
both branches. Singleton maps and the canonical coproduct reassociations
give the invertible compositor and unit comparison. Their coherence is
the monad unit/associativity equation on labels and equality of the unique
singleton maps. In particular, no nontrivial sample space was quotiented
to a singleton.

The embeddings agree on the designated pure maps, up to their canonical
unit identifications. The context-only and content-only specializations
obtained by replacing the opposite whole algebra with the identity algebra
are also recovered by these formulas. This specialization is distinct from
full faithfulness inside the fixed mixture. In particular $T=\id$ does
not trivialize receipt content: its fibres are still $\Set/B$.
\end{proof}

\subsection{The fixed monad history control}\label{pc:fixed-T}

Let $H^T_R(A,B)$ have objects $(P,h:E_P\to TB)$ and arrows
$\alpha:P\to P'$ satisfying $h=h'E_\alpha$. An endpoint square over
$(u,v)$ satisfies $Tv\,h=h'\bar\alpha$. This retains the full context
witness category, rather than just terminal output distributions.

For $(P,h):A\to B$ and $(R,k):B\to C$, write
$h(a,x)=(\sigma,(b_i))$. Its composite has context
$\coprod_x\prod_iR_{b_i}$ and body
\begin{equation}\label{pc:fixed-composition}
 (a,x,(r_i))\longmapsto\mu_C T(k)(\sigma,(b_i,r_i)).
\end{equation}
The unit has singleton context and body $\eta$. Composite cell functions
are the context functions in Section~\ref{pc:composition} with receipt
components omitted.

A three-stage history consists of a first sample, one second-stage sample
at every position of its $T$-output, and one third-stage sample at every
position of each resulting $T$-output. Cartesianness of $\mu$ identifies
these nested positions bijectively with the positions of the flattened
structure. Its associativity equation makes the two flattenings agree.
The canonical associator is precisely that history rebracketing. Its weight
is unchanged: both sides take the product over the same sampling nodes.
Units insert singleton nodes. The pentagon and triangle compare the same
addressed histories and consequently equal functions with equal weights.
This gives the fixed-$T$ retained-history composition with its canonical
associator. It is the terminal-content-grade $T$ control.

\begin{proposition}\label{pc:receipt-forgetting}
There is a pure-compatible strong comparison
\[
 E:H_{R,T}\longrightarrow H^T_R,\qquad
 E(P,Q,f)=(P,\nu_Qf),\quad E(\alpha,\beta)=\alpha.
\]
It preserves all endpoint-square operations. In general it is neither
locally full nor locally faithful.
\end{proposition}
\begin{proof}
For a source endpoint square,
$\nu_{Q'}(v+\beta)=Tv\,\nu_Q$. Thus the target cell satisfies exactly
$Tv\,\nu_Qf=\nu_{Q'}f'\bar\alpha$. Identity and vertical composition
are preserved.

The contexts of $E(F;G)$ and $EF;EG$ are the same: both sample the
positions of $\nu_Qf$. Equation~\eqref{pc:label-composite} identifies
their bodies with~\eqref{pc:fixed-composition}, since the later body is
$\nu_Sg$. Their compositor is the identity context cell, or its canonical
rebracketing representative with weak choices. It is invertible. The mixed
identity is sent to the singleton-$\eta$ identity, giving the unit comparison.
The explicit composite context cell of Section~\ref{pc:composition} is
also the target composite cell, proving naturality for all source cells
and endpoints.

Equivalently, the content-only map
$\psi(Q,f)=\nu_Qf$ is a morphism to the fixed-$T$ content algebra:
its unit equation is $\nu_{\varnothing}=\eta$, and its multiplication
equation is the label equation defining accumulated receipts. At a generic
crossing the equality
\[
 \nu_{\widehat Q}\chi(y,(x_i))=t\in T(E_R)
\]
identifies the receipt crossing followed by $\psi$ with the fixed-$T$
crossing of $\nu_Q(y)$. Both retain the same product context. This is the
compatibility square of the two law packages; its unit and multiplication
comparisons are the same addressed rebracketings. Thus the comparison also
agrees with the inherited functorial construction on compatible law maps.

For three arrows, either route keeps the same three-level sample tree and
performs the canonical rebracketing of Section~\ref{pc:fixed-T}. Receipt
labels are replaced by their $T$-labels; equation~\eqref{pc:label-composite}
turns receipt substitution into $T$-multiplication. Associativity of $\mu$
identifies the resulting labels. A target native cell is a context map,
and both routes give the identical addressed-history function. This proves
the strong comparison pentagon. Removing singleton samples and applying
the unit laws proves its triangle. Horizontal square pasting is the
composite context formula already checked, and vertical pasting is
composition of context maps. Pure maps are preserved. Hence the comparison
is strong at the realized and endpoint-square levels, not merely sound for
closed output laws.

For nonfaithfulness take $A=B=\one$, singleton contexts, receipt sets
$Q=\{q\}$ and $Q'=\{q_0,q_1\}$ all labelled $\eta_B(*)$, and bodies
returning the ordinary output $*$. The two maps $q\mapsto q_0$ and
$q\mapsto q_1$ are distinct native cells and have the same image under $E$.

For nonfullness take one arrow with empty receipt set and ordinary body
$*$, and another with one receipt $q$ labelled $\eta_B(*)$ whose body
returns $q$. Their images are the same fixed-$T$ history arrow. Its identity
cell has no source lift, since the sole map $\varnothing\to\{q\}$ cannot
turn an ordinary output into a receipt. These examples also exclude an
equivalence claim.
\end{proof}

Every fixed-$T$ arrow has an object-level lift by taking $Q=TB$, $q=\id$
and $f=\inr h$. This does not supply full faithfulness. For $T=\mathrm{List}$,
restricting to the full-support rational or nonnegative-real weighted
subcategory recovers the previously used probability/List control.
Products and dependent accumulation preserve full support there. With
zero-weight witnesses retained, the target is its raw-weight extension,
not literally the earlier restricted presentation. The comparison is via
$\nu$, not a claim that fixing a List grade in a container category is
closed under substitution.

These controls do not replace the complete mixed-law proof. They specify
no change of $T$, no QBS representation, no physical implementation, and
no application semantics. Receipt forgetting retains sample histories but
forgets receipt cells; coefficient change retains those functions while
forgetting numerical information. An additional kernel quotient would be a
third, separate comparison.

\section{A genuinely dependent indexed example}\label{ex:indexed}

Take $R=\mathbb Q_{\ge0}$ and $T=\List$. Let the input interface be
$A=\one$ and $B=\{b_0,b_1\}$. A first fair-bit context has sample $x$.
Its receipt set is $Q=\{a_0,a_1,c\}$, with
\[
 q(a_0)=q(a_1)=[b_0,b_1],\qquad q(c)=[b_1].
\]
Its body returns $a_0$ if $x=0$ and $c$ if $x=1$. The receipt $a_1$ is
unused by this body but remains part of its grade. Let the continuation
context at $b_0$ be a fair bit and at $b_1$ a uniform four-point witness.
The continuation may have empty receipt grade and a constant ordinary
output into $C=\one$.

The composite context is a dependent sum. The $x=0$ branch has eight
points $(0,r_0,r_1)$, each of weight $1/16$; the $x=1$ branch has four
points $(1,r_1)$, each of weight $1/8$. Thus no fixed product independent
of the first result describes the continuation. Its full lifted receipt
grade has eight decorations over each of $a_0,a_1$ and four over $c$:
twenty tags, including the eight unused alternatives over $a_1$.

Now apply the native receipt map
\[
 a_0,a_1\longmapsto m,\qquad c\longmapsto n,
 \qquad q'(m)=[b_0,b_1],\quad q'(n)=[b_1],
\]
and the context map that is the identity at $b_0$ and parity
$\{0,1,2,3\}\to\{0,1\}$ at $b_1$. These are both noninvertible
native cells. Their composite context function is
\[
 (0,r_0,r_1)\longmapsto(0,r_0,r_1\bmod2),\qquad
 (1,r_1)\longmapsto(1,r_1\bmod2).
\]
Each binary target history has two preimages of mass $1/16$, hence mass
$1/8$; each unary target history has two preimages of mass $1/8$, hence
mass $1/4$. There are four binary and two unary target histories, with
total mass one. The receipt map sends twenty decorated tags to six,
preserving their complete List labels. No probability measure is assigned
to the receipt set: weight preservation is a condition on the context
fibres, not a spurious sum over all tags.

This example uses result-dependent arity, an unused receipt alternative,
receipt merging, and a noninvertible context map in the same computation.
It does not import affine position deletion into slice cells: the binary
tags map to a binary tag, and the unary tag maps to a unary tag. A pure
identification of $b_0$ and $b_1$ would change labels but still retain two
distinct sampling occurrences in the binary branch.

\section{What the theorem establishes and what it does not}\label{sec:verdict}

The independent input test is concrete: give a commutative semiring with
$0\ne1$, raw finite normalized witnesses, a finite-position polynomial
presentation of $T$, and cartesian unit and multiplication. It is applicable
before constructing a mixed composition. It yields the full law, its
retained-cell relative algebra, and the comparisons proved above; no law
or equivalent already-coherent lift has been included among the inputs.

The sufficient criterion and bounded necessity have a precise common
interface. The canonical occurrence-display and tag-preserving reassignment
recipe detects cartesianness. The stronger converse applies to the declared
occurrence-product, raw-witness, $T(E)$-label-retaining normal-form class;
it is not recognition of arbitrary unmarked bicategories or an impossibility
result for every alternative interaction. The prefix-two example separates
the ordinary output-distribution law from invertible structural comparisons
with retained witnesses. It is not merely another restatement of the supplied
mixed-law factory.

\begin{longtable}{@{}p{0.28\linewidth}p{0.64\linewidth}@{}}
\toprule
Commitment & Established strength or exact boundary \\
\midrule
\endhead
Independent positive family &
The semiring and cartesian polynomial monad are specified independently.
The crossing, its native action, pure balance, endpoints and complete
C1--C10 coherence are derived. \\
Genuine indexing and dependence &
$L_B=\Set/(TB)$ varies with $B$. Continuations depend on prior samples
and their polynomial outputs; Section~\ref{ex:indexed} gives a finite
example with both native grades active. \\
Bounded recognition &
The canonical content recipe has an exact cartesianness test; the fixed
retained normal-form class has the stronger unit/multiplication necessity.
Changing that interface can invalidate the converse. \\
Nonidentity input comparison &
Every unital $R\to R'$ at fixed $T$ induces a pure-compatible strong
comparison. Positive rational weights to Boolean possibilities is
nonidentity and generally not locally full. No $T$-changing result is claimed. \\
Conservative one-sided controls &
Both unit inclusions are strong and locally fully faithful, including
their endpoint squares. Recovery when the opposite whole algebra is
trivialized is proved separately. \\
Inherited probability/List control &
The fixed-$T$ map $(P,Q,f)\mapsto(P,\nu_Qf)$ is strong but generally
neither full nor faithful. For List, restricting to full-support weights
recovers the previous presentation; raw weights extend it. \\
Earlier display/codisplay branch &
The pullback/pushout programme remains an open comparison and obstruction
branch. It is not solved by the present polynomial receipt theorem. \\
Full expanded research contract &
No interpreted O6 application family, no-recovery semantics, language,
QBS comparison, external monoidal action or complete integrated contract
is established here. Boolean coefficients alone do not supply them. \\
\bottomrule
\end{longtable}

The known ordinary mechanism is credited to the primary-source baseline,
not advertised as new. The retained indexed lift, its exact normal-form
boundary and its controls are the mathematical claims proved in this
dossier. This is a successful proof of the authorized TARGET B package,
not an assertion that every obligation in O1--O8 has been executed or that
the full research contract can already be accepted.

\section{Verification and reproducibility}\label{sec:verification}

The complete crossing proof, preservation arguments and bounded recognition
were independently reviewed adversarially at the stated scope. The review
found no unresolved mathematical gap in this theorem package. This is a
human-readable mathematical proof and review, not proof-assistant
formalization or an exhaustive computational verification of the universal
statements.

The accompanying standard-library Python checks use exact finite arithmetic.
They test the prefix-two monad on a two-element base, including all 1,806
triple-nested elements; selected ordinary traversal unit, multiplication and
naturality equations; the sixteen-to-four retained projection; and the
twelve-to-six dependent context map with the twenty-to-six full receipt map
in Section~\ref{ex:indexed}. A separate retained-associativity checker tests
50 context histories and 131 receipt tags for a triple, then 273 context
histories and 761 receipt tags for a full pentagon, including 141 zero-weight
histories and 490 receipt alternatives unreachable from the body.

These finite checks are regression evidence, not substitutes for the
parameterized proofs. The source checker verifies balanced TeX grouping and
environments, references, citations and the absence of assembly placeholders.
The TeX is supplied uncompiled, as permitted; no PDF or visual-layout verification
is asserted. The reproducibility package contains the final TeX, checks,
results and a manifest, not intermediate proof debates or primary-source PDFs.



\begin{thebibliography}{10}
\bibitem{Contract}
\emph{Expanded Research Contract Version 0.1}.
Project TeX source, 1 November 2026; provisionally adopted after drafting.
O1 is the independent compatibility and bounded recognition obligation.

\bibitem{Feasibility}
\emph{O1 Feasibility and Scope Decision}.
Project TeX source, 1 November 2026. TARGET B and its assumption ledger.

\bibitem{Unified}
\emph{One relative construction for context content and effects}.
Project TeX source, 1 November 2026.
Relevant labels: \texttt{thm:coend}, \texttt{def:law},
\texttt{thm:mixed}, \texttt{thm:problaw}.

\bibitem{FTMH}
S. Fujii, Y. C. Tsai, Y. Montacute and I. Hasuo.
\emph{Monads and Distributive Laws in Substructural Contexts}, extended
version, 2026. Theorem 42; Propositions 35--36; Examples 5 and 38.
\url{https://arxiv.org/html/2605.13533v1}.

\bibitem{GambinoLobbia}
N. Gambino and G. Lobbia.
\emph{On the formal theory of pseudomonads and pseudodistributive laws}.
Theory and Applications of Categories 37 (2021), no. 2, 14--56.
Version consulted: arXiv:0907.1359v3, 20 January 2021,
Definition 4.1, Remark 4.2 and Table 1, and Appendix A.
\url{https://arxiv.org/abs/0907.1359v3}.

\bibitem{Marmolejo}
F. Marmolejo.
\emph{Distributive laws for pseudomonads}.
Theory and Applications of Categories 5 (1999), no. 5, 91--147.
Section 4 and Theorem 5.1. The composite-pseudomonad theorem is used through
the inherited framework~\cite{Unified}; the law-equation correspondence is
checked by \cite[Remark 4.2 and Table 1]{GambinoLobbia}.
\end{thebibliography}
\end{document}
